% =============================================================================
% Results (RQ1-RQ3) + Discussion + Conclusion  -- drafted 2026-07-25
% Paste into the main .tex, replacing the empty RQ subsections and the missing
% Discussion/Conclusion.
%
% NUMBER PROVENANCE:
%   [V] = recomputed 2026-07-25 directly from results/P1_v4/{probe,behavior}_by_t.json
%         (balanced, n-weighted pooling over t) and from P3_postcot/responses/*.json
%         re-parsed with the current ANSWER_RE. Safe to quote.
%   [U] = carried from the 2026-07-23 session log, NOT yet re-verified against
%         disk. Marked inline as VERIFY. Do not submit these without rerunning.
%   [P] = pending BoS data (job 2894711 / the bos-pre-* sweep).
%
% IPD is excluded from the scenario suite per the 2026-07-25 decision: its
% optimal action is horizon- rather than belief-driven, so it does not measure
% the same construct as the other scenarios. See the Limitations subsection.
% =============================================================================

\section{Results}

We organize the results around three questions: whether models act rationally
at all (RQ1), whether their evidence tracking improves as history accumulates
(RQ2), and, where they fail, which stage of the decision process produced the
failure (RQ3). Unless stated otherwise, all numbers are balanced over the two
baseline-optimal action classes, pooled over decision steps, and computed on
unambiguous steps only ($|p_{\mathrm{BASELINE}}-0.5|>0.10$).

\subsection{RQ1: Do LLMs Make Rational Decisions?}

Table~\ref{tab:main-results} reports agreement with the baseline-optimal action
before any generation (Pre-CoT, P2). Two patterns dominate.

\textbf{\textit{Decision quality is governed by scenario structure, not by
scale alone.}}
On the pure inference scenarios, the largest models are close to the baseline:
Gemma-4 31B and Llama-3.3 70B both reach $100$ on \textit{Bandit} and $98$ and
$92$ on \textit{Monty Hall}, while the smallest models sit near chance on
\textit{Monty Hall} ($52$ for Gemma-2 2B) despite being far above chance on
\textit{Bandit} ($92$).\,[V] Scale therefore resolves pure inference. It does
not resolve the reactive-opponent scenarios: on \textit{Stag Hunt} no model
exceeds $79$, and the three largest models cluster at $76$, $76$, and $79$ --
a ceiling roughly $20$ points below their own performance on the inference
scenarios.\,[V]

\textbf{\textit{The deficit is not explained by the difficulty of the decision
rule.}}
A natural alternative account is that models fail wherever the decision
boundary is a nontrivial expected-utility threshold rather than a posterior
probability of $0.5$. The data do not support it. Grouping scenarios by
threshold rather than by structure gives differences that are inconsistent in
sign across models (threshold $1/2$ minus threshold $2/3$: $+13$, $+4$, $+9$,
$-13$, $+10$, $+4$, $+10$, $-3$, $-5$ points).\,[V] Models that recover the
$2/3$ threshold in \textit{Trust Game} and \textit{Machine} still fail in
\textit{Stag Hunt}, where the same threshold applies.

\subsection{RQ2: Do LLMs Update on Evidence Like Rational Agents?}

We measure evidence tracking as the Spearman rank correlation, within a
session, between the baseline posterior trajectory and the model's preference
trajectory over the ten decision steps. Rank correlation is the appropriate
statistic here because the normative claim is monotone tracking rather than
linearity, and it is invariant to any monotone rescaling of the model's logit
range.

Pre-CoT tracking is positive but uneven across families (Bayesian-inference
scenarios $\rho=0.54$, $n=3600$ sessions; game-theoretic scenarios
$\rho=0.29$, $n=1800$).\,[V] Models therefore do track accumulating evidence,
and the tracking is weaker exactly where decision quality is weaker. The
per-scenario ordering is \textit{Bandit} $0.75$, \textit{Trust Game} $0.51$,
\textit{Monty Hall} $0.47$, \textit{Machine} $0.45$, \textit{Chicken} $0.30$,
\textit{Stag Hunt} $0.28$.\,[V]

The more consequential finding concerns what reasoning does to this tracking.
Chain-of-thought \emph{degrades} evidence tracking in the Bayesian-inference
scenarios ($\rho=0.54\rightarrow0.31$) while leaving it essentially unchanged
in the game-theoretic ones ($0.29\rightarrow0.26$).\,[V] The mechanism is
saturation: post-CoT logit gaps sit near ceiling regardless of how strong the
evidence actually is. Combined with RQ3 below, this yields a dissociation worth
stating plainly: \emph{self-generated reasoning improves accuracy while making
expressed confidence less informative}. Accuracy and calibration move in
opposite directions, which complicates the common assumption that more
reasoning yields a better-calibrated decision maker.
% CAVEAT to state in the text or a footnote: the post-CoT correlation is
% computed only over sessions with at least three parseable answers, so models
% that frequently fail to terminate are under-represented in that pool
% (n=1267 and n=511 sessions vs 3600 and 1800 pre-CoT).

Confidence magnitude grows with accumulated evidence only for the three largest
models. Averaged over the six scenarios, the mean absolute pre-CoT logit gap
between step~2 and step~10 rises by $+1.42$ (Qwen-3.6 27B), $+2.12$ (Gemma-4
31B) and $+1.82$ (Llama-3.3 70B) nats, while every sub-10B model is flat
($+0.04$ to $+0.24$) and R1-Qwen-7B declines ($-0.35$).\,[V] Accumulating
evidence therefore sharpens the internal preference only once models are large
enough, which is the same capability boundary that separates the models that
map beliefs to actions from those that do not.

\subsection{RQ3: When Decisions Are Wrong, Where Do They Fail?}

\textbf{\textit{Beliefs are present at every scale.}}
The belief probe (P1) decodes the baseline-optimal action from the residual
stream at $98$--$100$ balanced accuracy on every Bayesian-inference scenario,
for every model including the 2B anchor, and at $86$--$97$ on \textit{Stag
Hunt}.\,[V] Belief formation is, on this suite, essentially solved. It is not
where models differ.

\textbf{\textit{The failure is localized to the belief-to-action mapping.}}
Comparing P1 with P2 within the same forward pass isolates the mapping stage.
For the mid-sized and large models the belief-to-action drop is $4$ points on
the Bayesian-inference scenarios ($100\rightarrow96$) and $17$ points on the
reactive-opponent and anti-coordination scenarios ($97\rightarrow80$).\,[V]
The information required for the correct action is linearly present in the
representation and is not used by the mechanism that selects the action. The
gap is directionally consistent: \emph{all nine models} show a larger
belief-to-action loss in the game-theoretic scenarios than in the
Bayesian-inference ones, by $7$, $16$, $9$, $24$, $16$, $16$, $9$, $21$, and
$31$ points respectively.\,[V]

\textbf{\textit{The deficit is not social framing.}}
Because the game-theoretic scenarios are also the socially framed ones, an
obvious confound is that instruction tuning installs social priors that
override the computed belief. The \textit{Trust Game}--\textit{Machine} pair
was constructed to test exactly this: matched hidden types, matched utilities,
and outcome trajectories paired by seed, differing only in whether the
counterparty is described as a person or a machine. On identical evidence the
three largest models differ by $+5$, $+1$, and $+5$ points between the social
and mechanical framings.\,[V] Narrative sociality alone produces no
systematic effect. What distinguishes the failing scenarios is that the
environment responds to the agent's own previous action, so the decision-relevant
quantity is a posterior predictive over the opponent's next move rather than a
posterior over a static parameter.

\textbf{\textit{Reasoning partially repairs the mapping.}}
Post-CoT agreement (P3) exceeds pre-CoT agreement in most cells, and the repair
is largest where the pre-CoT mapping was worst: Qwen-3.6 27B moves from $76$ to
$98$ on \textit{Stag Hunt} and from $83$ to $100$ on \textit{Chicken}.\,[V]
The repair is not uniform, however: Llama-3.3 70B \emph{declines} from $79$ to
$71$ on \textit{Stag Hunt} while improving from $69$ to $98$ on
\textit{Chicken}.\,[V] Reasoning is therefore best described as a second,
partially independent attempt at the mapping rather than as a reliable
correction to it.

\textbf{\textit{Reasoning-distilled models fail by not terminating.}}
The two R1-distilled models frequently fail to produce any parseable final
answer, and the rate at which they terminate tracks their own pre-CoT mapping
competence rather than scenario family. R1-Qwen-7B terminates on $51\%$ of
\textit{Bandit} steps but only $2$--$6\%$ of steps in the reactive-opponent
scenarios; scaling the distillation to 70B rescues termination on the
inference scenarios ($91\%$ on \textit{Bandit}) while the reactive-opponent
scenarios remain largely non-terminating.\,[V] When the belief-to-action
mapping is broken, chain-of-thought does not act as an independent remedy: it
has nothing coherent to converge on, and manifests as unproductive looping
rather than as corrected output.

\section{Discussion}

\textbf{\textit{Localization changes the remedy.}}
The practical value of separating the stages is that they imply different
interventions. Because beliefs are linearly decodable at every scale we test,
prompting strategies aimed at improving inference address a stage that is not
failing. The stage that fails is the mapping from a represented belief to a
selected action, and it fails hardest where the correct action depends on a
posterior predictive over a responsive environment. Interventions that
externalize the decision rule -- supplying the posterior explicitly, or
scaffolding the expected-utility comparison -- target the stage the measurements
implicate; interventions that ask the model to ``think more carefully about the
evidence'' do not.

\textbf{\textit{Why the reactive scenarios are harder.}}
In the Bayesian-inference scenarios the optimal action is a monotone function
of a posterior over a static hidden parameter. In the reactive scenarios the
agent must additionally condition on its own last action, because the
opponent's next move depends on it. The belief probe shows that models
represent the relevant state; what they appear not to do reliably is carry that
state through the additional conditioning step that the payoff comparison
requires. This is consistent with the failure surviving both scale and social
de-framing.

\textbf{\textit{Accuracy is not calibration.}}
The RQ2 and RQ3 results point in opposite directions and should be reported
together. Chain-of-thought raises agreement with the optimal action while
flattening the relationship between expressed confidence and evidence strength.
For deployed agents that gate actions on model-reported confidence, the
implication is direct: reasoning traces improve the decision and degrade the
signal one would use to decide whether to trust it.

\subsection{Limitations}

\textbf{\textit{Committed opponents.}}
Opponents follow fixed strategies rather than best-responding. The
game-theoretic scenarios are therefore Bayesian decision problems against
committed types rather than equilibrium computations. This is what makes an
exact optimal action available at every step, and it means our claims concern
best response to a predicted opponent, not equilibrium selection.

\textbf{\textit{A scenario we designed out.}}
An earlier version of this suite included the Iterated Prisoner's Dilemma. We
excluded it after establishing that its optimal action is not belief-driven:
because defection strictly dominates in the stage game under the payoffs used
($T=5$, $R=3$, $P=1$, $S=0$), the expected-utility difference
$\mathrm{EU}(D)-\mathrm{EU}(C)=p_c+1\ge1$ for every posterior, so no belief can
overturn it. Restoring belief relevance requires a finite-horizon continuation
term, which makes the label a function of the remaining horizon as well as the
posterior: the same posterior yields opposite optimal actions at different
steps, and a predictor that sees only the step index and never the history
reproduces the resulting labels on $70.5\%$ of steps, against a base rate of
$56.3\%$.\,[V] The corresponding clock-only gain is $+0.0$ points for
\textit{Stag Hunt} and $+2.5$ points for \textit{Bach or Stravinsky}. Since
the scenario does not measure the same construct as the rest of the suite, we
report it separately rather than in the main table. We replaced it with
\textit{Bach or Stravinsky}, which retains a reactive opponent while
guaranteeing structurally that the label cannot depend on the horizon: with no
dominant action, the posterior predictive changes the argmax within the round.

\textbf{\textit{Scope.}}
All scenarios use a binary action space and a fixed horizon of ten steps.
Histories are generated by a fixed behavior policy, so our claims concern
observer-style decision making: whether a model responds rationally to
provided evidence, not whether its own policy would have generated informative
evidence.

\section{Conclusion}

We introduced a pipeline that localizes where rational decision making breaks
down inside an LLM, by pairing sequential decision scenarios that admit exactly
computable Bayesian-optimal actions with three logit-level checkpoints aligned
to the stages of the reference procedure. Across nine models from 2B to 70B,
belief formation is essentially solved: the optimal action is linearly
decodable from the residual stream at $98$--$100$ accuracy at every scale.
Decision quality nonetheless varies by more than fifty points across scenarios,
and the variation is concentrated in a single stage -- the mapping from the
represented belief to the selected action -- which degrades sharply when the
environment responds to the agent's own previous action, and which is not
explained by model scale, by the difficulty of the decision rule, or by social
framing. Chain-of-thought partially repairs this mapping while simultaneously
making expressed confidence less informative about the evidence. Progress on
LLM agents in interactive settings, on this evidence, depends less on better
inference than on connecting inference the models already perform to the
actions they take.
